\documentclass{article}
\usepackage{graphicx} % Required for inserting images
\usepackage{xcolor}
\usepackage[draft]{changes}
% \usepackage[colorinlistoftodos]{todonotes}
\usepackage{todonotes}
\usepackage{natbib}
\usepackage{amsmath}
\usepackage{amsfonts}
\usepackage{hyperref}
\usepackage{makecell}
\usepackage[table]{xcolor}
\usepackage{caption}
\usepackage{booktabs}
\usepackage{float}


\title{Project Description}
\author{Anand Gnanadesikan, Somdatta Goswami, Jan Drgona}
\date{January 2026}

\begin{document}

\section{Introduction}

The ocean features many phenomena which emerge from the interaction of processes across space and time scales. One example is the global overturning circulation, which emerges from the interaction between surface heat fluxes, atmospheric freshwater transport, mixing by internal waves and mesoscale eddies, and surface wind. Another is  El Nino, which emerges from the coupling between waves in the atmosphere and ocean, but which is modulated by a multitude of small-scale processes. Yet another is the emergence of biomes from the interaction of large-scale physical circulation, seasonal cycling of light, and the regeneration of nutrients within the ocean interior. 

A critical goal of oceanography is to develop parsimonious yet accurate theories that describe key aspects of these emergent phenomena. A classic paper by \citep{stommel1961thermohaline} reduced the overturning circulation to two closed-form quadratic equations in which the transport of water vapor in the atmosphere controlled the overturning in the ocean and which exhibited tipping point behavior. \cite{gnanadesikan1999simple} \deleted{Gnanadesikan (1999)} extended this theory to include the role of Southern Ocean winds, mesoscale eddies and mixing, describing the circulation in terms of a cubic equation. Similarly, it is also possible to reduce El Nino to a simple set of coupling equations between the ocean and atmosphere \citep{vallis1986nino} \deleted{(Vallis, 1986; von Oldenburgh et al., 2005)}, with key variables being the strength of equatorial winds, temperature gradient along the equator and depth of the oceanic pycnocline.  Allometric (size-structured) theories for the scaling between environmental variables and biomass have been shown to capture large-scale patterns in the biomass of phytoplankton, which make up the base of the marine food web (Dunne et al., 2005; Holder and Gnanadesikan, 2023). However, such theories often involve poorly constrained or arbitrarily chosen parameters and may oversimply the systems they describe. {\bf The goal of this proposal is to develop a methodology for using modern ML methods to discover more robust equation sets from large datasets and to apply them to oceanographic questions.}  

Reducing oceanographic systems to more parsimonious sets of equations advances our understanding of the ocean in several ways. First, the process of doing so can identify key control parameters. For example, a key insight from \cite{gnanadesikan1999simple} was that the northward transport and warming of surface waters by winds within the Southern Ocean acted to increase the overturning circulation because it resulted in the transformation of dense waters to light waters. Because this transport scales as the wind stress $\tau_s$  it focussed attention on whether the Earth System Models used to simulate the global climate system properly captured this wind stress. As shown in subsequent work, many did not \citep{russell2006southern}, thus allowing us to use the theory to understand why models differ. This can be particularly important in cases where the responses to changes are non-intuitive. For example, \cite{gnanadesikan2017impact} found that increasing lateral mixing in a coupled model increased the amplitude of warm El Nino anomalies along the equator despite mixing acting to damp such anomalies. The reason was that a net warming of the eastern Pacific produced an atmosphere that was more responsive to such anomalies.  

Finally, identifying emergent equations allows the identification of manifolds in parameter space along which different control parameters may give similar results. \cite{gnanadesikan1999simple} noted that the northward transport of light water could be opposed by the tendency of meoscale eddies (the equivalent of atmosphere storms in the ocean) to flatten tilted density surfaces. Following \cite{mcwilliams1990linear}\deleted{Gent and McWilliams (1990)} this tendency was parameterized as being proportional to a thickness diffusion coefficient $A_{GM}$. However, this in turn meant that one could match the observed overturning and density structure with different combinations of $\tau_s$ and $A_{GM}$. Simply because a solution happened to match a particular set of observations did not mean it was correct. Identification of such tradeoffs enabled targetted experiments to evaluate what aspects of physical or biogeochemical solutions lying along such a manifold in parameter space were robust to these changes and which ones were not. For example \cite{mignone2006central}\deleted{Mignone et al. (2006)} found that while the total uptake of carbon dioxide by the ocean (a process that slows global warming by about 25\%) was not sensitive to compensating changes in 
$\tau_s$ and  $A_{GM}$  
, the distribution of air-sea carbon flux (something that might be important to monitor) was in fact sensitive. The \cite{gnanadesikan1999simple} paper also showed that a second pathway of transformation of dense water into light waters associated with downward mixing of heat in the tropics could also play a role in driving the overturning, though a lesser role than previously assumed. The current generation of climate models still shows a range regarding which pathways dominate, with models dominated by the low-latitude upwelling pathway showing more of a tendency to have the North Atlantic overturning collapse under global warming \citep{baker2024wind}\deleted{(Baker et al., 2024)}. 

A challenge in developing reduced equation sets, however, is that a number of somewhat arbitrary choices can go into developing them. In \cite{gnanadesikan1999simple} the assumption was made that the wind stresses at the latitude at the tip of South America were critical- but subsequent work \citep{allison2010winds}\deleted{(Allison, Johnson and Marshall, 2010)} showed that what really mattered were winds over the northern edge of the Circumpolar Current. While the difference did not matter for the original experimental design in which the winds were scaled by a constant factor within the Southern Ocean (reflecting uncertainties in satellite recoveries of wind stress at the time) it remains relevant for understanding historical changes in which winds shifted Southward as a result of the ozone hole. Similarly, it is possible to develop models of El Nino by fitting equations that produce a delayed oscillator, but the results depend on the delay that one uses \citep{russell2006southern}\deleted{(Russell and Gnanadesikan, 2014)}. With respect to biogeochemical modeling \cite{dunne2005empirical}\deleted{Dunne et al. (2005)} found it possible to fit multiple allometric equations for phytoplankton using the relatively small ($\mathbb{O}(\sim 100)$) dataset of in-situ measurements available at the time. 

Machine learning methods offer a new way in which to develop such emergent equations. This proposal brings together three investigators at Johns Hopkins who have been working to do this in different ways. Below we describe recent work which we plan to bring together as part of this proposal. 

As already noted, lead PI Gnanadesikan has worked on developing emergent equations for both physical and biogeochemical phenomena in the oceans, including the three use cases on which we will focus here.  Co-PI Somdatta Goswami's group at Johns Hopkins has focused on equation discovery for canonical dynamical systems using neural-operator frameworks that decouple functional-form identification from parameter estimation. Finally, co-I Jan Drgona's group has Jan Drgona's group at JHU, currently leads development of the  Neuromancer SciML library \citep{drgona2023domain} in PyTorch for solving constrained optimization, physics-informed machine learning, and optimal control problems.

This work has been based on two key insights. The first, that the growth rate $\mu$ of phytoplankton is strongly dependent on light, nutrients, and temperature has been understood for decades \citep{riley1949volume}, though it has not been obvious how to scale widely divergent results from single species in a lab can be scaled up to ecosystem levels. The second is that the relative change in biomass from month to month is much smaller than the growth rate, so that the net loss rate 
$$
\lambda(P)\approx \mu
$$
 
Insofar as  $\lambda(P)$  is invertible, one might expect to be able to link observed biomass to environmental conditions. In a recent series of papers, the Gnanadesikan group has shown that it is possible to use ML emulators to do this for both models \citep{holder2021can, holder2022using, holder2023well} and satellite-based observations \citep{holder2023well, jin2024using, dutta2025using}  

\section{AI for Oceans}
Research into applications of artificial intelligence (AI) and machine learning (ML) in ocean, atmospheric, and climate sciences has accelerated in recent years \citep{eyring2024pushing, 10505123}. The applications of ML can be broadly classified into categories (A) classification and anomaly detection, which is concerned with, e.g., finding extreme event patterns or the classification of important structures or regimes; (B) regression, which is concerned with state reconstruction of important state variables, parameters, or diagnostics (metrics) from available data; and (C) state prediction, ranging from nowcasting to operational forecasting and sub-seasonal to seasonal prediction. A comprehensive collection of review articles on deep learning in Earth sciences is given in Camps et al., 2021 \cite{camps2021deep}, covering algorithmic foundations and examples of all three categories.

Classification in oceanography involves assigning labels to data like image of a fish, a sound from a whale, or a specific type of sea ice. Deep learning models (like CNNs) are used to identify fish species, coral health (bleaching vs. healthy), and plankton types from underwater ROV footage. Acoustic applications like distinguishing between different whale species or even specific pods by analyzing hydrophone data. Benthic Habitat Mapping is another application where sonar backscatter is combined with depth data to classify the seafloor into categories like sandy, rocky reef, or seagrass meadow. Satellite SAR (Synthetic Aperture Radar) data is used for classifying sea ice types (first-year vs. multi-year ice), which is critical for safe Arctic navigation and climate modeling. Identifying physical structures like eddies, fronts, and upwelling zones from satellite sea-surface temperature (SST) maps is another example of ML application. Anomaly detection is another area which has witnessed application of ML. These applications focuses on finding outliers data points that deviate significantly from the norm, often signaling environmental threats or rare events. Some examples are 
environmental monitoring like detecting Pollution \& Oil Spills, plastic patches from  satellite imagery that differ from surrounding water, Harmful Algal Blooms by detecting unusual spikes in chlorophyll or water temperature, Marine Heatwaves using time-series analysis (like LSTMs), Maritime Security (Dark Vessels) using AIS (Automatic Identification System) data to find anomalous behavior, such as a ship turning off its transponder or loitering in a protected marine area—often a sign of illegal fishing. The architectures that have used are typically CNNs, LSTMs/ RNNs, Autoencoders, U-Net. 

State reconstruction of important state variables, parameters, or diagnostics from available data is another application area in ocean sciences where ML has been very prevalent. The most cutting-edge research in ocean state reconstruction has moved beyond simple statistical regression into Scientific Machine Learning (SciML). A recent study has used StrAss-PINNs (Stratified Assimilation PINNs) \citep{limousin2025deep} to reconstruct ocean flows in both the interior and surface layers, demonstrating a strong ability to resolve key ocean mesoscale features, including vortex rings, eastward jets associated with potential vorticity fronts, and smoother Rossby waves.
This allows to reconstruct high-resolution flow in the "abyss" using only sparse data from Argo floats and surface satellites. Another example is the use of temporal convolutional neural network (TCN) to predict 4D temperature fields (3D space + time) \citep{champenois2024machine}. Another study uses SpadeUp, a novel deep learning model that fuses surface data (wind fields, sea surface height, and surface currents) with subsurface thermohaline data to achieve high-precision 3D ocean current reconstruction \citep{li2025deep}.

Operational ocean prediction follows similar framework as that of numerical weather prediction (NWP). Its core engine is a data assimilation (DA) framework, consisting of a physical model (i.e., a complex algorithm for solving a set of partial differential equations, PDEs), a workflow for quality-controlling and ingesting diverse observational data streams into the DA system (ideally in near-real time), and an optimal estimation algorithm that combines models and data in a formal manner that produces statistically optimal forecasts. ML approaches can be incorporated into accelerating various elements of the DA workflow. Opportunities in the application of ML for partial differential equation (PDE) based models fall into two main categories, one being targeted insertion of ML within a physical model, and the other being complete replacement of the physical model by a surrogate model. In the former, certain elements or sub-components of a physical model are replaced by a surrogate model (e.g., a neural network), whereas in the latter, the entire model is emulated. A major source of model uncertainty is the parameterization of subgrid-scale (SGS) processes in terms of structural errors (formulation of functional representations of parameterizations) and parametric uncertainties (calibrating empirical parameters in the functional representations). Replacement of conventional functional representation subgrid-scale (SGS) turbulent oceanic processes with surrogate models that are based on machine learning is one of the frontline applications of ML in oceanography \citep{zanna2025framework, bodner2025data, sane2023parameterizing}. Similarly, equation discovery to infer the functional form of such SGS ocean parameterization schemes is another area of active study \citep{zanna2020data, perezhogin2023implementation}. A longer list of related efforts exists for numerical weather prediction \citep{ben2024rise, dueben2021deep}. These surrogates, mostly some form of neural networks, have been trained on (i.e., fit to) what are considered simulations of much higher fidelity and where these processes are resolved (e.g., large eddy simulations). Related efforts aim at learning improved parameterizations from online bias correction or analysis increments incurred in sequential data assimilation \citep{gregory2023deep, gregory2024machine, storto2024correction}. The second application of hybrid approaches is the to replace specific numerical algorithms within PDE based models with surrogate models to accelerate the simulation's time to solution. Such studies already exist in the fields of computational fluid dynamics \citep{kochkov2021machine} and atmospheric modelling \citep{arcomano2023hybrid}. These methods take advantage of differentiable programming to generate code for derivative of the physics model particularly the adjoint model that enable efficient real-time learning of the model parameters. Embedding of ML approaches within operational data assimilation workflows deployed in ocean prediction is another area of active research. 

Fully data driven approaches has been widely popular in the field of weather forecasting. The most popular approaches have been use of CNNs, Graph Neural Networks, Transformers, Fourier neural operators, RNNs/ LSTMs. Probabilistic approaches using variational autoencoders, GANs and diffusion models have also been employed. PINNs have also being used though their performance for extrapolation has generally been sub-optimal. The application of ML in ocean state forecasting has been much less as compared to that of weather forecasting. Some prominent examples are using LSTMs, RNNs, hybrid PINN-LSTM  \citep{minuzzi2023deep, lawal2025optimizing} for surface wave prediction,  surface wave–current interaction forecasting, storm surge forecasting \citep{xie2023developing}, and sea surface temperature prediction via deep learning \citep{wolff2020statistical, xu2023deep} and the use of neural networks for accelerating resonant nonlinear wave–wave interaction in an ocean surface wave model, regional to coastal sea level prediction, ocean color mapping, and statistical downscaling. Other applications include estimating ocean surface circulation \citep{sinha2021estimating, subel2024building}.  Recent studies have used  Fourier neural operators for modeling oceanic processes \citep{bire2023ocean, chattopadhyay2024oceannet, sun2024streamlining}. These studies present fully data-driven ocean models that
match the capabilities of traditional numerical ocean models in predicting high-resolution sea surface height (SSH) fields.  Transformer based applications have also been used recently for ocean eddy-resolving forecasting,  sea ice cover
in the polar oceans \citep{wang2024xihe, andersson2021seasonal}. Challenges in the domain of ocean state prediction using Ml has been seen in extending the applications to seasonal, inter-annual, and multi-decadal timescales. Here, the increased need for models or invariant operators (physics-based or surrogates) to conserve fundamental properties (mass, energy, momentum, and active tracers) puts severe demands on ML approaches.

\section{Neural ODE Performance}
The availability of adequate time series observed data for environmental drivers is a challenge. Therefore, we resorted to using the Global Ocean Physics Reanalyses (GLORYS12V1) monthly dataset from the Copernicus Marine Service. The chlorophyll-a data was obtained from the Global Ocean Colour Plankton L4 dataset which contains monthly average mass concentration of chlorophyll-a. Data between 2015-2022 was taken for the analysis. To determine that the data has adequate information to model the plankton dynamics and also to establish a reference performance criteria, a Neural ODE (NODE) architecture was used for modeling the dynamics as a first step. The NODE for plankton dynamics can be written as below.
    \begin{align*}
    \frac{\partial P_i}{\partial t} = f(P_i, Dr, Nut;\theta)
    \label{eq:plankton_node}
    \end{align*}
Here, $P_i$ represents the concentration or abundance of the $i^{th}$ species of plankton (chlorophyll-a in our case); $Dr$ accounts for environmental drivers, such as physical conditions like sea surface temperature, mixed layer depth, light etc.; $Nut$ denotes the availability of nutrients like nitrate, phosphate, silicate; and $f$ is the neural network parametrized by $\theta$. The NODE architecture learns a continuous time differential equation instead of learning a a sequence of discrete layers. NODE learns the plankton dynamics. Although being a black box model, the NODE gives valuable understanding of dynamical systems and acts as a baseline for further more advanced analysis. The details of the NODE architecture training is given in Table-\ref{tab:node_trg}. We found that NODE performed reasonably well as can be understood from the performance metrics in Table-\ref{tab:node_trg}. 


\begin{table}[h!]
    \centering
    \captionsetup{font={large,bf}}
    \caption{NODE Training}
    \rowcolors{2}{gray!10}{white}
    \resizebox{\columnwidth}{!}{
    \begin{tabular}{@{}>{\bfseries}p{3.2cm} p{7.8cm}@{}}
    \toprule
    Parameter & Value \\ 
    \midrule
    Total Samples      & 9868 \\
    Train Samples      & 6567 \\
    Test Samples       & 3,301 \\
    Time Span          & 0–60 months \\
    Time Steps         & 1 month \\
    Species            & Chlorophyll-a \\
    Integrator         & Runge–Kutta 4 (RK4) \\
    Architecture       & 3 hidden layers, 256 units each \\
    Activation         & SiLU \\
    Learning Rate      & Adaptive Learning Rate \\
    Optimizer          & Adam (with weight decay) \\
    Batch Size         & 32 \\
    Curriculum Stages  & [12, 24, 48] months \\
    Epochs per Stage   & [20, 30, 40, 100] \\
    Metrics            & RMSE = 0.0563, $R^2$ = 0.7380, Pearson = 0.8577 \\
    Loss Function      & MSE + Onestep + FD + Init + Derivative \\
    Dropout Rate       & 0.2 \\
    Training Method    & BPTT \\
    \bottomrule
    \end{tabular}
    }
    \label{tab:node_trg}
\end{table}

\section{SINDy}
Sparse Identification of Nonlinear Dynamics (SINDy) is a data-driven technique used to identify non-linear dynamical systems. SINDy operates under the presumption that a sparse set of basis functions can adequately characterize the dynamics of a system. Sparse regression methods can then be used to calculate the coefficients of these basis functions. The final model will be described in terms of a set of sparse ordinary differential equations that define the dynamics of the system \cite{pandey2024data}. Various factors such as the system’s complexity, data, and selection of the hyperparameters, affect how well SINDy captures transient and steady-state behaviour. Unlike "black-box" nonparametric techniques like neural networks, which may require no assumptions but often lack interpretability, SINDy produces parsimonious models—sparse ordinary differential equations (ODEs) that balance accuracy with model complexity to avoid overfitting. The core assumption of SINDy is that the system dynamics, represented as
\[
\frac{d}{dt} x(t) = f(x(t))
\]
are sparse in the space of possible functions. This means that while there are infinitely many ways to describe a system's motion, the "true" physical laws typically involve only a handful of active terms \cite{brunton2016discovering}. By leveraging advances in compressed sensing and machine learning, SINDy identifies these active terms from a high-dimensional library of candidate functions without requiring computationally expensive brute-force searches. SINDy uses optimization algorithms—such as Sequential Thresholded Least Squares (STLSQ), Stepwise Sparse Regression (SSR), or Sparse Relaxed Regularized Regression (SR3)—to determine which terms in the library are non-zero. This process effectively "trims" the library to retain only the functions that significantly impact the system's dynamics. SINDy has demonstrated a remarkable ability to capture both fast transient dynamics and slow steady-state behaviors, provided the training data covers both regimes. The effectiveness of SINDy is highly dependent on data quality and the completeness of the feature library. While the method is efficient with relatively small datasets, its performance declines as noise levels increase, which can lead to biased coefficients or overfitting. Furthermore, if the "true" model structure is not included in the feature library, SINDy may fail to isolate the correct dynamics.We plan to use SINDy or one of its extensions/ modifications like weak SINDy, constrained SINDy as the main technique for system identification in all the three problem statements.

\section{Drgona: Scientific machine learning} (1 page+figure) 
\textcolor{violet}{Dr. Drgoňa is also leading Johns Hopkins as part of the new multi-institution Learning Accelerated Domain Sciences (LEADS) Institute, a Department of Energy-funded collaborative initiative focused on reshaping how artificial intelligence supports scientific discovery. This initiative is supported through the Scientific Discovery through Advanced Computing (SciDAC) program which seeks to build new AI methods and software tools aimed at solving some of the nation’s most pressing scientific and engineering challenges, including those related to sustainable energy, infrastructure, and materials science. Dr. Drgoňa is a member of the Institute of Electrical and Electronics Engineers and the Association for Computing Machinery. He regularly serves as a reviewer for related journals including Applied Energy, Automatica, IEEE Control Systems Letters, IEEE Transactions on Control Systems Technology, IEEE Transactions on Industrial Informatics, Control Engineering Practice, Journal of Process Control, Energy and Buildings, Journal of Control Automation and Electrical Systems, and Electric Power Systems Research.}\todo[inline]{[SD]This is arbitrary text which may be used if and where found suitable}

\section{Goswami: DeepONet} (1 page+figure) 

Recent work in the Goswami group has focused on equation discovery for canonical dynamical systems using neural-operator frameworks that decouple functional-form identification from parameter estimation. Building on this foundation, we will employ a Deep Operator Network (DeepONet) architecture to uncover compact, interpretable dynamical equations linking atmospheric circulation patterns to ocean overturning and gyre variability in the North Atlantic. DeepONet provides a flexible framework for learning mappings between high-dimensional spatio-temporal fields, such as atmospheric forcing or sea-surface temperature anomalies and corresponding dynamical responses or modal amplitudes. Because these datasets are inherently high-dimensional and multi-scale, we will first apply reduced-order modeling techniques, e.g., autoencoder, to identify low-dimensional latent spaces that capture the dominant structures across different spatial and temporal scales. Within these latent manifolds, we will deploy a deep hidden-physics model to infer distinct sets of governing equations characterizing dynamics at each scale. By integrating physics-aware regularization with sparsity-promoting inversion, the learned operator enables (i) hidden-physics discovery, via symbolic or sparse regression on the surrogate to extract governing terms, and (ii) parameter identification, through efficient inverse estimation of closure coefficients across model resolutions (see Goswami et al., 2025; Kag et al., 2025; Mandl et al., 2025). The resulting framework will yield compact parametric representations with quantified uncertainties that can be systematically compared across resolutions and validated for dynamical fidelity before re-integration into coupled climate model. 

\textcolor{violet}{Dr. Goswami's research interests lies on advancing the fields of Scientific Computing, Computational Mechanics, and Machine Learning, encompassing fundamental and applied aspects. Her  group, Centrum IntelliPhysics, focuses on Scientific Machine Learning, and Artificial Intelligence for Science, contributing to developing algorithms and architectures tailored to address computational challenges in engineering, physical, and biological systems. She is interested in developing methods to address long-time horizon problems and challenges of coupling scales in multiscale multiphysics material modeling. The group conducts research on developing AI-accelerated numerical simulations to enhance the efficiency and accuracy of these processes. Dr. Goswami’s research has received funding from the Department of Energy. She received the DARPA INTACT Grant in 2025, the National Science Foundation’s National Artificial Intelligence Research Resource (NAIRR) Pilot, and the Johns Hopkins University Discovery Award 2024.} \todo[inline]{[SD]This is arbitrary text which may be used if and where found suitable}

\section{Proposed work:}

 

\section{Use cases:}

 

\subsection{Ocean biogeochemical models:}  

Key goal- sets of equations that match observations better that can deployed in Earth System Models. Large amount of data available, understand the variables involved. 

\todo[inline]{Show Sandupal’s preliminary results} 

As already noted, the ocean overturning circulation plays a critical role in Earth’s climate, particularly within the North Atlantic.  In this basin, relatively warm surface waters are transported northward by the Gulf Stream and the North Atlantic Current into subpolar regions beyond roughly $45^\circ N$.  \textcolor{violet}{(red lines in Fig. 1a)}. During winter, these waters encounter cold, dry air masses originating from North America, Greenland, and Europe, releasing substantial heat to the atmosphere. This ocean–atmosphere exchange helps maintain the Northern Hemisphere’s warmth relative to the Southern Hemisphere at comparable latitudes and contributes to the striking temperature contrast between northwestern Europe and northeastern North America. As the surface waters cool and become denser, they sink into the deep ocean and flow southward \textcolor{violet}{(blue lines in Fig. 1a)}, driving the Atlantic Meridional Overturning Circulation (AMOC), a critical component of the global climate system that ventilates the deep ocean with oxygen and sequesters anthropogenic carbon. However, current generations of climate models disagree on the AMOC’s mean strength, period and amplitude of variability, and sensitivity to external forcing. The fundamental causes of these discrepancies remain elusive, reflecting uncertainties in how interacting physical processes spanning atmosphere, ocean, and cryosphere are represented across scales. 

AMOC variability involves coupling between the ocean and atmosphere. The release of heat associated with the overturning circulation helps drives a region of high storminess over the central Atlantic and an intense low pressure system known as the Icelandic Low (brown lines in Fig. 1a). The winds flowing around this low drive the counterclockwise oceanic circulation known as the subpolar gyre, of which the cold Labrador current makes up the western boundary. Analyses of the variation of the overturning often link variation in the overturning to differences in the strength of this low relative to the high pressure system centered over the Azores. But this “North Atlantic Oscillation” is only one of many ways in which the Icelandic Low can vary. Changes in the shape of the low, either extending towards the east or south, also have the potential to change ocean circulation and heat transport.  Gnanadesikan et al. (2020) found that Principal Oscillation Pattern analysis, in which the time series of trends in the amplitudes of rotated Empirical Orthogonal Functions are regressed against those amplitudes to derive oscillatory patterns, were able to recover a 50 year oscillation in overturning in open waters in the Southern Ocean within a particular model. However, we have since found that this method does not do as well in the North Atlantic, in part because the multitude of modes of variability tends to bias regression coefficients in this method to low values, unrealistically lengthening the predicted period of variation.  

We believe  

 

\subsection{Low-latitude climate variability} 

 

 

\subsection{Joint class-scientific machine learning and equation discovery for the environmental sciences} 

 

\section{Results of previous work} 

{\bf Anand Gnanadesikan}: Grant OCE-1756808 (6/2018-5/2022) Collaborative Research: Southern Ocean Convection in climate models: controls and Impacts (Lead PI Irina Marinov, 387K to JHU)  

{\bf Intellectual merit}

The goal of this proposal was to identify the dynamics underlying convective variability in the Southern Ocean in coupled climate models, to understand why these vary across different climate models, and to understand the impacts of such variability. A key result (published in Gnanadesikan et al., 2020a)  was the description of a coupled mode of variability in which: 
1. Convection in the central Weddell Sea warms the atmosphere, resulting in winds shifting northward. This allows freshwater to build up in the Eastern Weddell Sea as less freshwater is advectively exported and less salty water is mixed up from below. 2. The  resulting freshwater anomaly propagates into the Central Weddell Sea and shuts off convection.  3. This cools the Antarctic and shifts the winds southwards. 4. Which then results in a salty anomaly developing in the Eastern Weddell that propagates into the Central Weddell and turns the convection back on.

We showed that this mode was very sensitive to the representation of eddy mixing. High values of  ARedi  reduce vertical gradients, prevent salinity anomalies from developing and keep convection in the “on” state, High values of AGM increase stratification in the Central Weddell to the point where convection is always “off”-salinity anomalies never get big enough to break through. We also showed that the mean states of these models differed in how it represented the Antarctic Circumpolar Current \citep{ragen2020impact}\deleted{(Ragen et al, 2020a)}, and the sensitivity and linearity of hypoxia to global warming \citep{bahl2020scaling}\deleted{(Bahl et al., 2019a, 2020a)}. This work contributed to a review article summarizing the disconnect between models and observations of isopycnal mixing and why it might matters both for convection and for a broader range of processes \citep{abernathey2022deep}\deleted{(Abernathey et al., 2022)}. We also used these experiments to show that machine learning methods could distinguish intermodal differences in ocean biogeochemical cycling caused by changes in circulation from those caused by differences in the formulation of the biological model \citep{holder2022using}\deleted{(Holder et al., 2022a)}.

{\bf Data availability}  
New datasets/code generated as part of four publications \citep{gnanadesikan2020feedbacks, ragen2020impact, bahl2018understanding, bahl2020scaling}\deleted{(Gnanadesikan et al., 2020b, Ragen et al. 2022b, Bahl et al., 2018, 2020b)} have been made available on the Johns Hopkins Research Data archive and one dataset was made available via Zenodo \citep{holder2021can}\deleted{(Holder et al., 2021)}. 

{\bf Broader impacts}

The grant helped support masters student Alexis Bahl (who after working for a few years is now finishing up a Ph.D. in biological oceanography at UBC) and  Ph.D student Chris Holder at JHU. Additionally, four students who participated in this work as undergraduates ended up as authors or coauthors on publications. 

 

\bibliographystyle{ieeetr}
\bibliography{project.bib}
\end{document}